% =============================================================================
% FORMATTED LATEX TABLE IV AND PUBLICATION ANALYSIS FOR MAIN.TEX (SECTION V-D1)
% =============================================================================

\begin{table}[t]
\centering
\caption{Competing Risks Ablation proving naive filtering of planning failures creates safety bias distortion.}
\label{tab:competing_risks_ablation}
\footnotesize
\begin{tabularx}{\columnwidth}{@{}X c c c@{}}
\toprule
\textbf{Formulation Metric} & \textbf{Naive Filtering (Dropped $B^H=1$)} & \textbf{Competing Risks (Ours)} & \textbf{Bias Distortion} \\
\midrule
Evaluated Failure Rate & 1.66\% & 1.84\% & +9.7\% \\
\bottomrule
\end{tabularx}
\end{table}

To evaluate the mathematical validity of our hazard estimator, we conduct an ablation comparing naive conditional filtering against our proposed competing-risks formulation. Naively dropping probe steps that resulted in path planning failure ($B^H=1$) measures a collision probability of only $1.66\%$. However, total mission failure incorporating both collisions ($Y^H=1$) and path blockages ($B^H=1$) is $1.84\%$. Ignoring competing path failures creates a $9.7\%$ false-safety bias distortion, as hazardous trajectories that were prevented from executing by local planner collapse are erroneously treated as absent risk. Our competing-risks formulation eliminates this post-treatment collider bias, guaranteeing unbiased safety predictions.
